\section{MQuAKE: A Retrieval-Dominant Propagation Boundary}

MQuAKE tests whether updated facts propagate correctly to direct edits, single-hop consequences, and multi-hop consequences \citep{zhong2023mquake}. In our frozen official runs, this regime is also retrieval-dominant.

\begin{table}[h]
\centering
\small
\begin{tabular}{llrrrrr}
\toprule
Model & Policy & Quality & Edit & Single-hop & Multi-hop & Forgetting \\
\midrule
FLAN-T5 Base & Retrieval only & 0.551 & 0.214 & 1.000 & 0.214 & 0.00005 \\
FLAN-T5 Base & Full router & 0.284 & 0.188 & 0.431 & 0.168 & 0.00979 \\
FLAN-T5 Base & Calibrated router & 0.284 & 0.188 & 0.431 & 0.168 & 0.00979 \\
FLAN-T5 Base & Router without delayed utility & 0.338 & 0.156 & 0.643 & 0.093 & 0.01086 \\
FLAN-T5 Small & Retrieval only & 0.551 & 0.214 & 1.000 & 0.214 & 0.00005 \\
FLAN-T5 Small & Full router & 0.276 & 0.185 & 0.433 & 0.150 & 0.00226 \\
FLAN-T5 Small & Calibrated router & 0.276 & 0.185 & 0.433 & 0.150 & 0.00226 \\
FLAN-T5 Small & Router without delayed utility & 0.319 & 0.022 & 0.705 & 0.031 & 0.00329 \\
\bottomrule
\end{tabular}

\caption{Official MQuAKE results. Retrieval remains strongest on overall quality, direct edit success, and multi-hop consequences. Suppressing delayed utility improves on the full router but still does not overtake retrieval.}
\label{tab:mquake-main}
\end{table}

On both FLAN backbones, \texttt{always\_retrieve} reaches 0.551 overall quality. The full router reaches only 0.284 on FLAN-T5-Base and 0.276 on FLAN-T5-Small, and the calibrated router is identical. Suppressing delayed utility improves overall quality to 0.338 and 0.319, but that still remains far below retrieval.

The benchmark-specific metrics reinforce the same interpretation. Retrieval reaches edit success 0.214, single-hop accuracy 1.000, and multi-hop accuracy 0.214 on both FLAN backbones. The full router lowers all three. The \texttt{router\_no\_v\_future} variant increases single-hop accuracy to 0.643 on FLAN-T5-Base and 0.705 on FLAN-T5-Small, but it does so by sacrificing direct edit success and especially multi-hop accuracy, which falls to 0.093 and 0.031. Forgetting also rises relative to retrieval. Paired sign tests again show retrieval dominating with effectively zero \(p\)-values in both official FLAN bundles.

The MQuAKE lesson is therefore negative but precise: delayed utility does not transfer safely from stable recurrence to propagation-heavy edits. Even removing delayed utility entirely is only a partial repair, because the strongest safe policy in this regime is still retrieval-dominant.
